% GENERATED FILE. Edit canonical lessons, then run npm run generate:books.
% canonical-source-hash: fnv1a64:dd75c5a74a5c1c22
% canonical-lessons: GU-C405-udan, GU-C405-paylat, GU-C405-get, GU-C405-sutkes, GU-C405-cakravo

\chapter{Flights and Gates}
\label{ch:gu-flights-and-gates}
\hlchaptermodality{405}

\begin{chapteropening}
\textbf{By the end of this chapter:} \emph{I can find my gate and talk about the flight.}

Complete the last lesson of chapter 405: I can find my gate and talk about the flight.
\end{chapteropening}

\section[\texorpdfstring{\gu{ઉડાન} (uḍān)}{ઉડાન (uḍān)}]{\gu{ઉડાન} (uḍān) — a flight}

\label{lesson:GU-C405-udan}



\begin{quote}
Before the new one: say the Gujarati for an uphill climb, then the Gujarati for to jump over.
\end{quote}

\subsection*{You'll want to know: \gu{ઉડાન}}
\textbf{\gu{ઉડાન}} — \emph{uḍān} — \textquotedblleft{}a flight\textquotedblright{}.

\begin{quote}
\gu{અમારી} \gu{ઉડાન} \gu{આવતીકાલે} \gu{સવારે} \gu{છે}.
\end{quote}

\begin{quote}
\emph{amārī uḍān āvtīkāle savāre che.}
\end{quote}

\begin{quote}
Our flight is tomorrow morning.
\end{quote}

A feminine noun.

\begin{cousinweb}
\textbf{\gu{ઉડાન}} (\emph{uḍān}) is built on \gu{ઊડવું}, \textquotedblleft{}to fly\textquotedblright{}, a tadbhava of Sanskrit \emph{uḍḍīyatē}, \textquotedblleft{}flies up\textquotedblright{} (CDIAL 1697): the flying itself, and then the journey by air.
\end{cousinweb}

\begin{grammarlens}[title={What you've built}]
One more word for getting around.
\end{grammarlens}

\subsection*{Your turn}
\begin{itemize}
\raggedright
  \item \emph{Say it:} \emph{uḍān}
  \item \emph{Say it:} \emph{uḍān}, once more
  \item \emph{Say it:} \emph{amārī uḍān āvtīkāle savāre che.} — Our flight is tomorrow morning.
  \item \emph{Recall:} say the Gujarati for an uphill climb, then the Gujarati for to jump over, then say \emph{uḍān} again
\end{itemize}

\begin{tcolorbox}[breakable,colback=teal!4,colframe=teal!35!black,title={Before you move on}]
What does \gu{ઉડાન} mean? (\textquotedblleft{}A flight\textquotedblright{}.) Say it once more.
\end{tcolorbox}
\section[\texorpdfstring{\gu{પાયલટ} (pāylaṭ)}{પાયલટ (pāylaṭ)}]{\gu{પાયલટ} (pāylaṭ) — a pilot}

\label{lesson:GU-C405-paylat}



\begin{quote}
Before the new one: say the Gujarati for to jump over, then the Gujarati for a flight.
\end{quote}

\subsection*{You'll want to know: \gu{પાયલટ}}
\textbf{\gu{પાયલટ}} — \emph{pāylaṭ} — \textquotedblleft{}a pilot\textquotedblright{}.

\begin{quote}
\gu{પાયલટે} \gu{વિમાન} \gu{ધીમે} \gu{નીચે} \gu{ઉતાર્યું}.
\end{quote}

\begin{quote}
\emph{pāyalṭe vimān dhīme nīce utāryũ.}
\end{quote}

\begin{quote}
The pilot brought the plane down slowly.
\end{quote}

A masculine noun.

\begin{cousinweb}
\textbf{\gu{પાયલટ}} (\emph{pāylaṭ}) is borrowed from English \emph{pilot}.
\end{cousinweb}

\begin{grammarlens}[title={What you've built}]
One more word for getting around.
\end{grammarlens}

\subsection*{Your turn}
\begin{itemize}
\raggedright
  \item \emph{Say it:} \emph{pāylaṭ}
  \item \emph{Say it:} \emph{pāylaṭ}, once more
  \item \emph{Say it:} \emph{pāyalṭe vimān dhīme nīce utāryũ.} — The pilot brought the plane down slowly.
  \item \emph{Recall:} say the Gujarati for to jump over, then the Gujarati for a flight, then say \emph{pāylaṭ} again
\end{itemize}

\begin{tcolorbox}[breakable,colback=teal!4,colframe=teal!35!black,title={Before you move on}]
What does \gu{પાયલટ} mean? (\textquotedblleft{}A pilot\textquotedblright{}.) Say it once more.
\end{tcolorbox}
\section[\texorpdfstring{\gu{ગેટ} (geṭ)}{ગેટ (geṭ)}]{\gu{ગેટ} (geṭ) — an airport gate; a gateway}

\label{lesson:GU-C405-get}



\begin{quote}
Before the new one: say the Gujarati for a flight, then the Gujarati for a pilot.
\end{quote}

\subsection*{You'll want to know: \gu{ગેટ}}
\textbf{\gu{ગેટ}} — \emph{geṭ} — \textquotedblleft{}an airport gate; a gateway\textquotedblright{}.

\begin{quote}
\gu{અમારું} \gu{વિમાન} \gu{પાંચ} \gu{નંબરના} \gu{ગેટ} \gu{પરથી} \gu{ઊપડશે}.
\end{quote}

\begin{quote}
\emph{amārũ vimān pā̃c nambarnā geṭ parthī ūpaḍśe.}
\end{quote}

\begin{quote}
Our plane will leave from gate number five.
\end{quote}

A masculine noun.

\begin{cousinweb}
\textbf{\gu{ગેટ}} (\emph{geṭ}) is borrowed from English \emph{gate}, at an airport the numbered gate where you board.
\end{cousinweb}

\begin{grammarlens}[title={What you've built}]
One more word for getting around.
\end{grammarlens}

\subsection*{Your turn}
\begin{itemize}
\raggedright
  \item \emph{Say it:} \emph{geṭ}
  \item \emph{Say it:} \emph{geṭ}, once more
  \item \emph{Say it:} \emph{amārũ vimān pā̃c nambarnā geṭ parthī ūpaḍśe.} — Our plane will leave from gate number five.
  \item \emph{Recall:} say the Gujarati for a flight, then the Gujarati for a pilot, then say \emph{geṭ} again
\end{itemize}

\begin{tcolorbox}[breakable,colback=teal!4,colframe=teal!35!black,title={Before you move on}]
What does \gu{ગેટ} mean? (\textquotedblleft{}An airport gate; a gateway\textquotedblright{}.) Say it once more.
\end{tcolorbox}
\section[\texorpdfstring{\gu{સૂટકેસ} (sūṭkes)}{સૂટકેસ (sūṭkes)}]{\gu{સૂટકેસ} (sūṭkes) — a suitcase}

\label{lesson:GU-C405-sutkes}



\begin{quote}
Before the new one: say the Gujarati for a pilot, then the Gujarati for an airport gate.
\end{quote}

\subsection*{You'll want to know: \gu{સૂટકેસ}}
\textbf{\gu{સૂટકેસ}} — \emph{sūṭkes} — \textquotedblleft{}a suitcase\textquotedblright{}.

\begin{quote}
\gu{મારી} \gu{સૂટકેસ} \gu{ઘણી} \gu{ભારે} \gu{છે}.
\end{quote}

\begin{quote}
\emph{mārī sūṭkes ghaṇī bhāre che.}
\end{quote}

\begin{quote}
My suitcase is very heavy.
\end{quote}

A feminine noun.

\begin{cousinweb}
\textbf{\gu{સૂટકેસ}} (\emph{sūṭkes}) is borrowed from English \emph{suitcase}.
\end{cousinweb}

\begin{grammarlens}[title={What you've built}]
One more word for getting around.
\end{grammarlens}

\subsection*{Your turn}
\begin{itemize}
\raggedright
  \item \emph{Say it:} \emph{sūṭkes}
  \item \emph{Say it:} \emph{sūṭkes}, once more
  \item \emph{Say it:} \emph{mārī sūṭkes ghaṇī bhāre che.} — My suitcase is very heavy.
  \item \emph{Recall:} say the Gujarati for a pilot, then the Gujarati for an airport gate, then say \emph{sūṭkes} again
\end{itemize}

\begin{tcolorbox}[breakable,colback=teal!4,colframe=teal!35!black,title={Before you move on}]
What does \gu{સૂટકેસ} mean? (\textquotedblleft{}A suitcase\textquotedblright{}.) Say it once more.
\end{tcolorbox}
\section[\texorpdfstring{\gu{ચકરાવો} (cakrāvo)}{ચકરાવો (cakrāvo)}]{\gu{ચકરાવો} (cakrāvo) — a circle in the air; a detour}

\label{lesson:GU-C405-cakravo}



\begin{quote}
Before the new one: say the Gujarati for an airport gate, then the Gujarati for a suitcase.
\end{quote}

\subsection*{You'll want to know: \gu{ચકરાવો}}
\textbf{\gu{ચકરાવો}} — \emph{cakrāvo} — \textquotedblleft{}a circle in the air; a detour\textquotedblright{}.

\begin{quote}
\gu{વિમાને} \gu{શહેર} \gu{ઉપર} \gu{બે} \gu{ચકરાવા} \gu{લીધા} \gu{અને} \gu{પછી} \gu{નીચે} \gu{ઊતર્યું}.
\end{quote}

\begin{quote}
\emph{vimāne śaher upar be cakrāvā līdhā ane pachī nīce ūtaryũ.}
\end{quote}

\begin{quote}
The plane circled twice over the city and then came down.
\end{quote}

A masculine noun.

\begin{cousinweb}
\textbf{\gu{ચકરાવો}} (\emph{cakrāvo}) seems to be built on Sanskrit \emph{cakra}, \textquotedblleft{}a wheel\textquotedblright{} (Monier-Williams): a round that a plane, a bird or a traveller makes before coming down or going on.
\end{cousinweb}

\begin{grammarlens}[title={What you've built}]
One more word for getting around.
\end{grammarlens}

\subsection*{Your turn}
\begin{itemize}
\raggedright
  \item \emph{Say it:} \emph{cakrāvo}
  \item \emph{Say it:} \emph{cakrāvo}, once more
  \item \emph{Say it:} \emph{vimāne śaher upar be cakrāvā līdhā ane pachī nīce ūtaryũ.} — The plane circled twice over the city and then came down.
  \item \emph{Recall:} say the Gujarati for an airport gate, then the Gujarati for a suitcase, then say \emph{cakrāvo} again
\end{itemize}

\begin{tcolorbox}[breakable,colback=teal!4,colframe=teal!35!black,title={Before you move on}]
What does \gu{ચકરાવો} mean? (\textquotedblleft{}A circle in the air; a detour\textquotedblright{}.) Say it once more.
\end{tcolorbox}
